\section{Background}

\subsection{The Trajectory-Coupling Threat}

Papers~5--8 of the VulnPrio sequence~\cite{paper5temporal,paper6capacity,paper7online}
all share one acknowledged internal-validity threat. At every window of
every cell, fleet evolution is driven by HygienePrio-full's
top-$K$ selection. EPSS-only, HRS-only, CVSS-only, and Random are scored
against the resulting backlog --- a backlog whose composition was
determined by HygienePrio's selection preferences, not by their own.

Each paper documented the threat but bounded its claims to the
HygienePrio-driven trajectory. None of them measured what would happen
under a counterfactual trajectory.

\subsection{Why It Could Change the Headline Conclusion}

Paper~6~\cite{paper6capacity} reported that HygienePrio-full's W6 P@50
collapses at $K \in \{100, 200\}$ on the HygienePrio-driven trajectory
and attributed the collapse to selection that exhausts the high-HRS
tail. The
mechanism Paper~6 named is recursive: HygienePrio-full preferentially
selects high-HRS-host pairs, those hosts get patched, the upper tail of
the HRS distribution thins, and HygienePrio-full's discriminative
signal degrades.

If this mechanism is correct, then on a non-HygienePrio-driven
trajectory --- one whose selection policy does not drain the positive
set HygienePrio is evaluated against ---
HygienePrio's signal should \textit{not} collapse, because no
mechanism would be draining the HRS upper tail. The Paper~6 headline
H4-rejection ``HygienePrio collapses at high capacity'' would then
require a substantial qualifier: it collapses only when HygienePrio
itself is the policy doing the patching.

\subsection{Counterfactual Trajectories and Non-Identifiability}

A fully factorial ``each method drives its own trajectory'' evaluation
introduces a non-identifiability problem: when method A drives its own
trajectory and method B drives its own trajectory, A and B are scored
against \textit{different} evolving backlogs with \textit{different}
per-window positive sets. The number ``$A$'s mean P@50 at $w = 6$ on
$T_A$'' is not directly comparable to ``$B$'s mean P@50 at $w = 6$ on
$T_B$''. Papers~5--8 chose the HygienePrio-driven baseline specifically
to avoid this issue.

The present paper takes a different stance: we report the trajectory-stratified
P@50 grid explicitly --- one P@50 number per (driver, scorer, window) cell ---
and ask the operational questions in a way that respects the
non-identifiability rather than abstracting it away.
